\subsection{\label{sec.sign.det}Determinants}

Determinants were introduced by Leibniz in the 17th century, and quickly
became one of the most powerful tools in mathematics. They remained so until
the early 20th century. There is a 5-volume book by Thomas Muir \cite{Muir}
that merely summarizes the results found on determinants... until 1920.

Most of these old results are still interesting and nontrivial. The relative
role of determinants in mathematics has declined mainly because other parts of
mathematics have \textquotedblleft caught up\textquotedblright\ and have
produced easier ways to many of the places that were previously only
accessible through the study of determinants.

As with anything else, we will just present some of the most basic results and
methods related to determinants. For more, see \cite{MuiMet60}, \cite{Zeilbe},
\cite[Chapter 6]{detnotes}, \cite[Chapter I]{Prasolov}, \cite[Chapter
9]{BruRys91} and various other sources. A good introduction to the most
fundamental properties is \cite{Strick13}.

\begin{convention}
\label{conv.det.K}For the rest of Section \ref{sec.sign.det}, we fix a
commutative ring $K$. In most examples, $K$ will be $\mathbb{Z}$ or
$\mathbb{Q}$ or a polynomial ring.
\end{convention}

\begin{convention}
\label{conv.det.matrices}Let $n,m\in\mathbb{N}$. \medskip

\textbf{(a)} If $A$ is an $n\times m$-matrix, then $A_{i,j}$ shall mean the
$\left(  i,j\right)  $-th entry of $A$, that is, the entry of $A$ in row $i$
and column $j$. \medskip

\textbf{(b)} If $a_{i,j}$ is an element of $K$ for each $i\in\left[  n\right]
$ and each $j\in\left[  m\right]  $, then
\[
\left(  a_{i,j}\right)  _{1\leq i\leq n,\ 1\leq j\leq m}%
\]
shall denote the $n\times m$-matrix whose $\left(  i,j\right)  $-th entry is
$a_{i,j}$ for all $i\in\left[  n\right]  $ and $j\in\left[  m\right]  $.
Explicitly:%
\[
\left(  a_{i,j}\right)  _{1\leq i\leq n,\ 1\leq j\leq m}=\left(
\begin{array}
[c]{cccc}%
a_{1,1} & a_{1,2} & \cdots & a_{1,m}\\
a_{2,1} & a_{2,2} & \cdots & a_{2,m}\\
\vdots & \vdots & \ddots & \vdots\\
a_{n,1} & a_{n,2} & \cdots & a_{n,m}%
\end{array}
\right)  .
\]


Note that the letters \textquotedblleft$i$\textquotedblright\ and
\textquotedblleft$j$\textquotedblright\ in the notation \textquotedblleft%
$\left(  a_{i,j}\right)  _{1\leq i\leq n,\ 1\leq j\leq m}$\textquotedblright%
\ are not carved in stone. We could just as well use any other letters
instead, and write $\left(  a_{x,y}\right)  _{1\leq x\leq n,\ 1\leq y\leq m}$
or (somewhat misleadingly, but technically correctly) $\left(  a_{j,i}\right)
_{1\leq j\leq n,\ 1\leq i\leq m}$ for the exact same matrix. (However,
$\left(  a_{j,i}\right)  _{1\leq i\leq m,\ 1\leq j\leq n}$ is a different
matrix. Whichever index is mentioned first in the subscript after the closing
parenthesis is used to index rows; the other index is used to index columns.)
\medskip

\textbf{(c)} We let $K^{n\times m}$ denote the set of all $n\times m$-matrices
with entries in $K$. This is a $K$-module. If $n=m$, this is also a
$K$-algebra. \medskip

\textbf{(d)} Let $A\in K^{n\times m}$ be an $n\times m$-matrix. The
\emph{transpose} $A^{T}$ of $A$ is defined to be the $m\times n$-matrix whose
entries are given by%
\[
\left(  A^{T}\right)  _{i,j}=A_{j,i}\ \ \ \ \ \ \ \ \ \ \text{for all }%
i\in\left[  m\right]  \text{ and }j\in\left[  n\right]  .
\]

\end{convention}

\subsubsection{Definition}

There are several ways to define the determinant of a square matrix. The
following is the most direct one:

\begin{definition}
\label{def.det.det}Let $n\in\mathbb{N}$. Let $A\in K^{n\times n}$ be an
$n\times n$-matrix. The \emph{determinant} $\det A$ of $A$ is defined to be
the element%
\[
\sum_{\sigma\in S_{n}}\left(  -1\right)  ^{\sigma}\underbrace{A_{1,\sigma
\left(  1\right)  }A_{2,\sigma\left(  2\right)  }\cdots A_{n,\sigma\left(
n\right)  }}_{=\prod\limits_{i=1}^{n}A_{i,\sigma\left(  i\right)  }}%
\]
of $K$. Here, as before:

\begin{itemize}
\item we let $S_{n}$ denote the $n$-th symmetric group (i.e., the group of
permutations of $\left[  n\right]  =\left\{  1,2,\ldots,n\right\}  $);

\item we let $\left(  -1\right)  ^{\sigma}$ denote the sign of the permutation
$\sigma$ (as defined in Definition \ref{def.perm.sign}).
\end{itemize}
\end{definition}

\begin{example}
For $n=2$, we have%
\begin{align*}
\det A  &  =\sum_{\sigma\in S_{2}}\left(  -1\right)  ^{\sigma}A_{1,\sigma
\left(  1\right)  }A_{2,\sigma\left(  2\right)  }\\
&  =\underbrace{\left(  -1\right)  ^{\operatorname*{id}}}_{=1}%
\underbrace{A_{1,\operatorname*{id}\left(  1\right)  }}_{=A_{1,1}%
}\underbrace{A_{2,\operatorname*{id}\left(  2\right)  }}_{=A_{2,2}%
}+\underbrace{\left(  -1\right)  ^{s_{1}}}_{=-1}\underbrace{A_{1,s_{1}\left(
1\right)  }}_{=A_{1,2}}\underbrace{A_{2,s_{1}\left(  2\right)  }}_{=A_{2,1}}\\
&  \ \ \ \ \ \ \ \ \ \ \ \ \ \ \ \ \ \ \ \ \left(
\begin{array}
[c]{c}%
\text{since the two elements of }S_{2}\text{ are the identity}\\
\text{map }\operatorname*{id}\text{ and the simple transposition }%
s_{1}=t_{1,2}%
\end{array}
\right) \\
&  =A_{1,1}A_{2,2}-A_{1,2}A_{2,1}.
\end{align*}
Using less cumbersome notations, we can rewrite this as follows: For any
$a,b,a^{\prime},b^{\prime}\in K$, we have%
\[
\det\left(
\begin{array}
[c]{cc}%
a & b\\
a^{\prime} & b^{\prime}%
\end{array}
\right)  =ab^{\prime}-ba^{\prime}.
\]


Similarly, for $n=3$, we obtain%
\[
\det\left(
\begin{array}
[c]{ccc}%
a & b & c\\
a^{\prime} & b^{\prime} & c^{\prime}\\
a^{\prime\prime} & b^{\prime\prime} & c^{\prime\prime}%
\end{array}
\right)  =ab^{\prime}c^{\prime\prime}-ac^{\prime}b^{\prime\prime}-ba^{\prime
}c^{\prime\prime}+bc^{\prime}a^{\prime\prime}+ca^{\prime}b^{\prime\prime
}-cb^{\prime}a^{\prime\prime}.
\]
(The six addends on the right hand side here correspond to the six
permutations in $S_{3}$, which in one-line notation are $123$, $132$, $213$,
$231$, $312$, and $321$, respectively.)

Similarly, for $n=1$, we obtain that the determinant of the $1\times1$-matrix
$\left(
\begin{array}
[c]{c}%
a
\end{array}
\right)  $ is%
\[
\det\left(
\begin{array}
[c]{c}%
a
\end{array}
\right)  =a.
\]
Here, the \textquotedblleft$\left(
\begin{array}
[c]{c}%
a
\end{array}
\right)  $\textquotedblright\ on the left hand side is a $1\times1$-matrix.

Finally, for $n=0$, we obtain that the determinant of the $0\times0$-matrix
$\left(  {}\right)  $ (this is an empty matrix, with no rows and no columns)
is%
\[
\det\left(  {}\right)  =\left(  \text{empty product}\right)  =1.
\]

\end{example}

Some (particularly, older) texts use the notation $\left\vert A\right\vert $
instead of $\det A$ for the determinant of the matrix $A$.

The above definition of the determinant is purely combinatorial: it is an
alternating sum over the $n$-th symmetric group $S_{n}$. Typically, when
computing determinants, this definition is not in itself very useful (e.g.,
because $S_{n}$ gets rather large when $n$ is large). However, in some cases,
it suffices. Here are a few examples:

\begin{example}
\label{exa.det.hollow5x5}Let $a,b,c,d,e,\ldots,p$ be $16$ elements of $K$.
Prove that%
\[
\det\left(
\begin{array}
[c]{ccccc}%
a & b & c & d & e\\
p & 0 & 0 & 0 & f\\
o & 0 & 0 & 0 & g\\
n & 0 & 0 & 0 & h\\
m & l & k & j & i
\end{array}
\right)  =0.
\]
(The \textquotedblleft$o$\textquotedblright\ is a letter \textquotedblleft
oh\textquotedblright, not a zero. Not that it matters much...)
\end{example}

\begin{proof}
[Proof of Example \ref{exa.det.hollow5x5}.]Let $A$ be the $5\times5$-matrix
whose determinant we are trying to identify as $0$; thus, $A_{1,1}=a$ and
$A_{1,2}=b$ and $A_{3,2}=0$ and so on. Notice that%
\begin{equation}
A_{i,j}=0\ \ \ \ \ \ \ \ \ \ \text{whenever }i,j\in\left\{  2,3,4\right\}
\label{pf.exa.det.hollow5x5.0}%
\end{equation}
(since $A$ has a \textquotedblleft hollow core\textquotedblright, i.e., a
$3\times3$-square consisting entirely of zeroes in its middle). We must prove
that $\det A=0$.

Our definition of $\det A$ yields%
\begin{equation}
\det A=\sum_{\sigma\in S_{5}}\left(  -1\right)  ^{\sigma}\prod\limits_{i=1}%
^{5}A_{i,\sigma\left(  i\right)  }. \label{pf.exa.det.hollow5x5.detA=}%
\end{equation}


Now, I claim that each of the addends in the sum on the right hand side is
$0$. In other words, I claim that $\prod\limits_{i=1}^{5}A_{i,\sigma\left(
i\right)  }=0$ for each $\sigma\in S_{5}$.

To prove this, fix $\sigma\in S_{5}$. The three numbers $\sigma\left(
2\right)  ,\sigma\left(  3\right)  ,\sigma\left(  4\right)  $ are three
distinct elements of $\left[  5\right]  $ (distinct because $\sigma$ is
injective), so they cannot all belong to the $2$-element set $\left\{
1,5\right\}  $ (since there are no three distinct elements in a $2$-element
set). Hence, at least one of them must belong to the complement $\left\{
2,3,4\right\}  $ of this set. In other words, there exists some $i\in\left\{
2,3,4\right\}  $ such that $\sigma\left(  i\right)  \in\left\{  2,3,4\right\}
$. This $i$ must then satisfy $A_{i,\sigma\left(  i\right)  }=0$ (by
(\ref{pf.exa.det.hollow5x5.0}), applied to $j=\sigma\left(  i\right)  $).
Thus, we have shown that there exists some $i\in\left\{  2,3,4\right\}  $ such
that $A_{i,\sigma\left(  i\right)  }=0$.

This shows that at least one factor of the product $\prod\limits_{i=1}%
^{5}A_{i,\sigma\left(  i\right)  }$ is $0$. Thus, the entire product is $0$.

Forget that we fixed $\sigma$. We thus have proved that $\prod\limits_{i=1}%
^{5}A_{i,\sigma\left(  i\right)  }=0$ for each $\sigma\in S_{5}$. Hence,%
\[
\det A=\sum_{\sigma\in S_{5}}\left(  -1\right)  ^{\sigma}\underbrace{\prod
\limits_{i=1}^{5}A_{i,\sigma\left(  i\right)  }}_{=0}=0,
\]
and thus Example \ref{exa.det.hollow5x5} is proved.

[We note that there are various alternative proofs, e.g., using Laplace
expansion. Also, if $K$ is a field, you can argue that $\det A=0$ using rank
arguments. See \cite[Exercise 6.47 \textbf{(a)}]{detnotes} for a
generalization of Example \ref{exa.det.hollow5x5}.]
\end{proof}

\begin{example}
Let $n\in\mathbb{N}$, and let $x_{1},x_{2},\ldots,x_{n}\in K$ and $y_{1}%
,y_{2},\ldots,y_{n}\in K$. Compute%
\[
\det\left(  \left(  x_{i}y_{j}\right)  _{1\leq i\leq n,\ 1\leq j\leq
n}\right)  =\det\left(
\begin{array}
[c]{cccc}%
x_{1}y_{1} & x_{1}y_{2} & \cdots & x_{1}y_{n}\\
x_{2}y_{1} & x_{2}y_{2} & \cdots & x_{2}y_{n}\\
\vdots & \vdots & \ddots & \vdots\\
x_{n}y_{1} & x_{n}y_{2} & \cdots & x_{n}y_{n}%
\end{array}
\right)  .
\]

\end{example}

Let us experiment with small $n$'s:%
\begin{align*}
\det\left(  {}\right)   &  =1;\\
\det\left(
\begin{array}
[c]{c}%
x_{1}y_{1}%
\end{array}
\right)   &  =x_{1}y_{1};\\
\det\left(
\begin{array}
[c]{cc}%
x_{1}y_{1} & x_{1}y_{2}\\
x_{2}y_{1} & x_{2}y_{2}%
\end{array}
\right)   &  =0;\\
\det\left(
\begin{array}
[c]{ccc}%
x_{1}y_{1} & x_{1}y_{2} & x_{1}y_{3}\\
x_{2}y_{1} & x_{2}y_{2} & x_{2}y_{3}\\
x_{3}y_{1} & x_{3}y_{2} & x_{3}y_{3}%
\end{array}
\right)   &  =0.
\end{align*}
This makes us suspect the following:

\begin{proposition}
\label{prop.det.xiyj}Let $n\in\mathbb{N}$ be such that $n\geq2$. Let
$x_{1},x_{2},\ldots,x_{n}\in K$ and $y_{1},y_{2},\ldots,y_{n}\in K$. Then,%
\[
\det\left(  \left(  x_{i}y_{j}\right)  _{1\leq i\leq n,\ 1\leq j\leq
n}\right)  =0.
\]

\end{proposition}

\begin{proof}
[Proof of Proposition \ref{prop.det.xiyj} (sketched).]The definition of the
determinant yields%
\begin{align*}
\det\left(  \left(  x_{i}y_{j}\right)  _{1\leq i\leq n,\ 1\leq j\leq
n}\right)   &  =\sum_{\sigma\in S_{n}}\left(  -1\right)  ^{\sigma
}\underbrace{\left(  x_{1}y_{\sigma\left(  1\right)  }\right)  \left(
x_{2}y_{\sigma\left(  2\right)  }\right)  \cdots\left(  x_{n}y_{\sigma\left(
n\right)  }\right)  }_{=\left(  x_{1}x_{2}\cdots x_{n}\right)  \left(
y_{\sigma\left(  1\right)  }y_{\sigma\left(  2\right)  }\cdots y_{\sigma
\left(  n\right)  }\right)  }\\
&  =\sum_{\sigma\in S_{n}}\left(  -1\right)  ^{\sigma}\left(  x_{1}x_{2}\cdots
x_{n}\right)  \underbrace{\left(  y_{\sigma\left(  1\right)  }y_{\sigma\left(
2\right)  }\cdots y_{\sigma\left(  n\right)  }\right)  }_{\substack{=y_{1}%
y_{2}\cdots y_{n}\\\text{(since }\sigma\text{ is a bijection }\left[
n\right]  \rightarrow\left[  n\right]  \text{)}}}\\
&  =\sum_{\sigma\in S_{n}}\left(  -1\right)  ^{\sigma}\left(  x_{1}x_{2}\cdots
x_{n}\right)  \left(  y_{1}y_{2}\cdots y_{n}\right) \\
&  =\left(  x_{1}x_{2}\cdots x_{n}\right)  \left(  y_{1}y_{2}\cdots
y_{n}\right)  \underbrace{\sum_{\sigma\in S_{n}}\left(  -1\right)  ^{\sigma}%
}_{\substack{=0\\\text{(by (\ref{eq.cor.perm.num-even.sum-sign}))}}}=0.
\end{align*}
Thus, Proposition \ref{prop.det.xiyj} is proved.
\end{proof}

As a consequence of Proposition \ref{prop.det.xiyj} (applied to $x_{i}=x$ and
$y_{j}=1$), we see that%
\begin{equation}
\det\underbrace{\left(
\begin{array}
[c]{cccc}%
x & x & \cdots & x\\
x & x & \cdots & x\\
\vdots & \vdots & \ddots & \vdots\\
x & x & \cdots & x
\end{array}
\right)  }_{\text{an }n\times n\text{-matrix with }n\geq2}=0
\label{eq.prop.det.xiyj.cor-x}%
\end{equation}
for any $x\in K$. In other words, if all entries of a square matrix of size
$\geq2$ are equal, then the determinant of this matrix is $0$.

\begin{example}
Let $n\in\mathbb{N}$, and let $x_{1},x_{2},\ldots,x_{n}\in K$ and $y_{1}%
,y_{2},\ldots,y_{n}\in K$. Compute%
\[
\det\left(  \left(  x_{i}+y_{j}\right)  _{1\leq i\leq n,\ 1\leq j\leq
n}\right)  =\det\left(
\begin{array}
[c]{cccc}%
x_{1}+y_{1} & x_{1}+y_{2} & \cdots & x_{1}+y_{n}\\
x_{2}+y_{1} & x_{2}+y_{2} & \cdots & x_{2}+y_{n}\\
\vdots & \vdots & \ddots & \vdots\\
x_{n}+y_{1} & x_{n}+y_{2} & \cdots & x_{n}+y_{n}%
\end{array}
\right)  .
\]

\end{example}

Let us experiment with small $n$'s:%
\begin{align*}
\det\left(  {}\right)   &  =1;\\
\det\left(
\begin{array}
[c]{c}%
x_{1}+y_{1}%
\end{array}
\right)   &  =x_{1}+y_{1};\\
\det\left(
\begin{array}
[c]{cc}%
x_{1}+y_{1} & x_{1}+y_{2}\\
x_{2}+y_{1} & x_{2}+y_{2}%
\end{array}
\right)   &  =-\left(  x_{1}-x_{2}\right)  \left(  y_{1}-y_{2}\right)  ;\\
\det\left(
\begin{array}
[c]{ccc}%
x_{1}+y_{1} & x_{1}+y_{2} & x_{1}+y_{3}\\
x_{2}+y_{1} & x_{2}+y_{2} & x_{2}+y_{3}\\
x_{3}+y_{1} & x_{3}+y_{2} & x_{3}+y_{3}%
\end{array}
\right)   &  =0.
\end{align*}
So we suspect the following:

\begin{proposition}
\label{prop.det.xi+yj}Let $n\in\mathbb{N}$ be such that $n\geq3$. Let
$x_{1},x_{2},\ldots,x_{n}\in K$ and $y_{1},y_{2},\ldots,y_{n}\in K$. Then,%
\[
\det\left(  \left(  x_{i}+y_{j}\right)  _{1\leq i\leq n,\ 1\leq j\leq
n}\right)  =0.
\]

\end{proposition}

\begin{proof}
[Proof of Proposition \ref{prop.det.xi+yj} (sketched).](See \cite[Example
6.7]{detnotes} for more details.) The definition of the determinant yields%
\begin{align*}
&  \det\left(  \left(  x_{i}+y_{j}\right)  _{1\leq i\leq n,\ 1\leq j\leq
n}\right) \\
&  =\sum_{\sigma\in S_{n}}\left(  -1\right)  ^{\sigma}\underbrace{\left(
x_{1}+y_{\sigma\left(  1\right)  }\right)  \left(  x_{2}+y_{\sigma\left(
2\right)  }\right)  \cdots\left(  x_{n}+y_{\sigma\left(  n\right)  }\right)
}_{\substack{=\prod_{i=1}^{n}\left(  x_{i}+y_{\sigma\left(  i\right)
}\right)  =\sum\limits_{I\subseteq\left[  n\right]  }\left(  \prod
\limits_{i\in I}x_{i}\right)  \left(  \prod\limits_{i\in\left[  n\right]
\setminus I}y_{\sigma\left(  i\right)  }\right)  \\\text{(by
(\ref{eq.lem.prodrule.sum-ai-plus-bi.eq}), applied to }a_{i}=x_{i}\text{ and
}b_{i}=y_{\sigma\left(  i\right)  }\text{)}}}\\
&  =\sum_{\sigma\in S_{n}}\left(  -1\right)  ^{\sigma}\sum\limits_{I\subseteq
\left[  n\right]  }\left(  \prod\limits_{i\in I}x_{i}\right)  \left(
\prod\limits_{i\in\left[  n\right]  \setminus I}y_{\sigma\left(  i\right)
}\right) \\
&  =\sum\limits_{I\subseteq\left[  n\right]  }\ \ \sum_{\sigma\in S_{n}%
}\left(  -1\right)  ^{\sigma}\left(  \prod\limits_{i\in I}x_{i}\right)
\left(  \prod\limits_{i\in\left[  n\right]  \setminus I}y_{\sigma\left(
i\right)  }\right) \\
&  =\sum\limits_{I\subseteq\left[  n\right]  }\left(  \prod\limits_{i\in
I}x_{i}\right)  \sum_{\sigma\in S_{n}}\left(  -1\right)  ^{\sigma}%
\prod\limits_{i\in\left[  n\right]  \setminus I}y_{\sigma\left(  i\right)  }.
\end{align*}


Now, I claim that the inner sum is $0$ for each $I$. In other words, I claim
that%
\begin{equation}
\sum_{\sigma\in S_{n}}\left(  -1\right)  ^{\sigma}\prod\limits_{i\in\left[
n\right]  \setminus I}y_{\sigma\left(  i\right)  }%
=0\ \ \ \ \ \ \ \ \ \ \text{for each }I\subseteq\left[  n\right]  .
\label{pf.prop.det.xi+yj.innersum}%
\end{equation}
For instance, for $I=\left\{  1,2\right\}  $, this is claiming that%
\[
\sum_{\sigma\in S_{n}}\left(  -1\right)  ^{\sigma}y_{\sigma\left(  3\right)
}y_{\sigma\left(  4\right)  }\cdots y_{\sigma\left(  n\right)  }=0.
\]


[\textit{Proof of (\ref{pf.prop.det.xi+yj.innersum}):} Fix a subset $I$ of
$\left[  n\right]  $. We shall show that all addends in the sum $\sum
_{\sigma\in S_{n}}\left(  -1\right)  ^{\sigma}\prod\limits_{i\in\left[
n\right]  \setminus I}y_{\sigma\left(  i\right)  }$ cancel each other -- i.e.,
that for each addend in this sum, there is a different addend with the same
product of $y_{j}$'s but a different sign $\left(  -1\right)  ^{\sigma}$. To
achieve this, we need to pair up each $\sigma\in S_{n}$ with a different
permutation $\sigma^{\prime}=\sigma t_{u,v}\in S_{n}$ that satisfies
$\prod\limits_{i\in\left[  n\right]  \setminus I}y_{\sigma^{\prime}\left(
i\right)  }=\prod\limits_{i\in\left[  n\right]  \setminus I}y_{\sigma\left(
i\right)  }$ but $\left(  -1\right)  ^{\sigma^{\prime}}=-\left(  -1\right)
^{\sigma}$. (Indeed, this pairing will then produce the required
cancellations: the addend for each $\sigma$ will cancel the addend for the
corresponding $\sigma^{\prime}$. To be more rigorous, we are here applying
Lemma \ref{lem.sign.cancel2} to $\mathcal{A}=S_{n}$, $\mathcal{X}=S_{n}$ and
$\operatorname*{sign}\sigma=\left(  -1\right)  ^{\sigma}\prod\limits_{i\in
\left[  n\right]  \setminus I}y_{\sigma\left(  i\right)  }$ (of course, this
should not be confused for the notation $\operatorname*{sign}\sigma$ for
$\left(  -1\right)  ^{\sigma}$) and \newline$f=\left(  \text{the map }%
S_{n}\rightarrow S_{n}\text{ that sends each }\sigma\in S_{n}\text{ to the
corresponding }\sigma^{\prime}\right)  $.)

So let us construct our pairing. Indeed, from $I\subseteq\left[  n\right]  $,
we obtain $\left\vert I\right\vert +\left\vert \left[  n\right]  \setminus
I\right\vert =\left\vert \left[  n\right]  \right\vert =n\geq3$; hence, at
least one of the two sets $I$ and $\left[  n\right]  \setminus I$ has size
$>1$. In other words, must be in one of the following two cases:

\textit{Case 1:} We have $\left\vert I\right\vert >1$.

\textit{Case 2:} We have $\left\vert \left[  n\right]  \setminus I\right\vert
>1$.

Let us first consider Case 1. In this case, we have $\left\vert I\right\vert
>1$. Thus, $\left\vert I\right\vert \geq2$. Pick two distinct elements $u$ and
$v$ of $I$. (These exist, since $\left\vert I\right\vert \geq2$.) Now, for
each permutation $\sigma\in S_{n}$, we set $\sigma^{\prime}:=\sigma t_{u,v}\in
S_{n}$. Then, each $\sigma\in S_{n}$ satisfies $\sigma^{\prime\prime
}=\underbrace{\sigma^{\prime}}_{=\sigma t_{u,v}}t_{u,v}=\sigma
\underbrace{t_{u,v}t_{u,v}}_{=t_{u,v}^{2}=\operatorname*{id}}=\sigma$. Hence,
we can pair up each $\sigma\in S_{n}$ with $\sigma^{\prime}\in S_{n}$. Any two
permutations $\sigma$ and $\sigma^{\prime}$ that are paired with each other
have different signs (indeed, if $\sigma\in S_{n}$, then $\sigma^{\prime
}=\sigma t_{u,v}$ and thus $\left(  -1\right)  ^{\sigma^{\prime}}=\left(
-1\right)  ^{\sigma t_{u,v}}=\left(  -1\right)  ^{\sigma}\underbrace{\left(
-1\right)  ^{t_{u,v}}}_{=-1}=-\left(  -1\right)  ^{\sigma}$), but the
corresponding products $\prod\limits_{i\in\left[  n\right]  \setminus
I}y_{\sigma\left(  i\right)  }$ and $\prod\limits_{i\in\left[  n\right]
\setminus I}y_{\sigma^{\prime}\left(  i\right)  }$ are equal (indeed, the
permutations $\sigma$ and $\sigma^{\prime}=\sigma t_{u,v}$ differ only in
their values at $u$ and $v$, but neither of these two values appears in any of
our two products, since $u,v\in I$). This shows that our pairing has precisely
the properties we want: Each $\sigma\in S_{n}$ satisfies $\prod\limits_{i\in
\left[  n\right]  \setminus I}y_{\sigma^{\prime}\left(  i\right)  }%
=\prod\limits_{i\in\left[  n\right]  \setminus I}y_{\sigma\left(  i\right)  }$
but $\left(  -1\right)  ^{\sigma^{\prime}}=-\left(  -1\right)  ^{\sigma}$ (and
thus $\sigma^{\prime}\neq\sigma$). As explained above, this completes the
proof of (\ref{pf.prop.det.xi+yj.innersum}) in Case 1.

Let us now consider Case 2. In this case, we have $\left\vert \left[
n\right]  \setminus I\right\vert >1$. Thus, $\left\vert \left[  n\right]
\setminus I\right\vert \geq2$. Pick two distinct elements $u$ and $v$ of
$\left[  n\right]  \setminus I$. (These exist, since $\left\vert \left[
n\right]  \setminus I\right\vert \geq2$.) We now proceed just as we did in
Case 1: For each permutation $\sigma\in S_{n}$, we set $\sigma^{\prime
}:=\sigma t_{u,v}\in S_{n}$. Then, each $\sigma\in S_{n}$ satisfies
$\sigma^{\prime\prime}=\underbrace{\sigma^{\prime}}_{=\sigma t_{u,v}}%
t_{u,v}=\sigma\underbrace{t_{u,v}t_{u,v}}_{=t_{u,v}^{2}=\operatorname*{id}%
}=\sigma$. Hence, we can pair up each $\sigma\in S_{n}$ with $\sigma^{\prime
}\in S_{n}$. Any two permutations $\sigma$ and $\sigma^{\prime}$ that are
paired with each other have different signs (this can be seen as in Case 1),
but the corresponding products $\prod\limits_{i\in\left[  n\right]  \setminus
I}y_{\sigma\left(  i\right)  }$ and $\prod\limits_{i\in\left[  n\right]
\setminus I}y_{\sigma^{\prime}\left(  i\right)  }$ are equal (indeed, the
permutation $\sigma^{\prime}=\sigma t_{u,v}$ can be obtained from $\sigma$ by
swapping the values at $u$ and $v$ (so that $\sigma^{\prime}\left(  u\right)
=\sigma\left(  v\right)  $ and $\sigma^{\prime}\left(  v\right)
=\sigma\left(  u\right)  $); thus, the products $\prod\limits_{i\in\left[
n\right]  \setminus I}y_{\sigma\left(  i\right)  }$ and $\prod\limits_{i\in
\left[  n\right]  \setminus I}y_{\sigma^{\prime}\left(  i\right)  }$ differ
only in the order in which their factors $y_{\sigma\left(  u\right)  }$ and
$y_{\sigma\left(  v\right)  }$ appear in them\footnote{Both factors
$y_{\sigma\left(  u\right)  }$ and $y_{\sigma\left(  v\right)  }$ do indeed
appear in these products, since $u$ and $v$ belong to $\left[  n\right]
\setminus I$.}; hence, these products are equal, since $K$ is commutative).
Again, this shows that our pairing has the properties we want, and thus the
proof of (\ref{pf.prop.det.xi+yj.innersum}) is complete in Case 2.

Thus, (\ref{pf.prop.det.xi+yj.innersum}) is proved in both cases.]

Now, we can finish our computation of the original determinant:%
\[
\det\left(  \left(  x_{i}+y_{j}\right)  _{1\leq i\leq n,\ 1\leq j\leq
n}\right)  =\sum\limits_{I\subseteq\left[  n\right]  }\left(  \prod
\limits_{i\in I}x_{i}\right)  \underbrace{\sum_{\sigma\in S_{n}}\left(
-1\right)  ^{\sigma}\prod\limits_{i\in\left[  n\right]  \setminus I}%
y_{\sigma\left(  i\right)  }}_{\substack{=0\\\text{(by
(\ref{pf.prop.det.xi+yj.innersum}))}}}=0.
\]
This proves Proposition \ref{prop.det.xi+yj}.
\end{proof}

\subsubsection{Basic properties}

The examples above show that pedestrian proofs of facts about determinants
(using just the definition) are possible, but become cumbersome fairly
quickly. Fortunately, determinants have a lot of properties that, once proved,
provide new methods for computing determinants.

Let us first recall a few basic facts that should (ideally) be known from a
good course on linear algebra. Recall that the transpose of a matrix $A$ is
denoted by $A^{T}$.

\begin{theorem}
[Transposes preserve determinants]\label{thm.det.transp}Let $n\in\mathbb{N}$.
If $A\in K^{n\times n}$ is any $n\times n$-matrix, then $\det\left(
A^{T}\right)  =\det A$.
\end{theorem}

\begin{proof}
See \cite[Corollary B.16]{Strick13} or \cite[Exercise 6.4]{detnotes} or
\cite[\S 5.3.2]{Laue-det}.
\end{proof}

\begin{theorem}
[Determinants of triangular matrices]\label{thm.det.triang}Let $n\in
\mathbb{N}$. Let $A\in K^{n\times n}$ be a triangular (i.e., lower-triangular
or upper-triangular) $n\times n$-matrix. Then, the determinant of the matrix
$A$ is the product of its diagonal entries. That is,%
\[
\det A=A_{1,1}A_{2,2}\cdots A_{n,n}.
\]

\end{theorem}

\begin{proof}
See \cite[Proposition B.11]{Strick13} or \cite[Exercise 6.3 and the paragraph
after Exercise 6.4]{detnotes}.
\end{proof}

As a consequence of Theorem \ref{thm.det.triang}, we see that the determinant
of a diagonal matrix is the product of its diagonal entries (since any
diagonal matrix is triangular).

\begin{theorem}
[Row operation properties]\label{thm.det.rowop}Let $n\in\mathbb{N}$. Let $A\in
K^{n\times n}$ be an $n\times n$-matrix. Then: \medskip

\textbf{(a)} If we swap two rows of $A$, then $\det A$ gets multiplied by
$-1$. \medskip

\textbf{(b)} If $A$ has a zero row (i.e., a row that consists entirely of
zeroes), then $\det A=0$. \medskip

\textbf{(c)} If $A$ has two equal rows, then $\det A=0$. \medskip

\textbf{(d)} Let $\lambda\in K$. If we multiply a row of $A$ by $\lambda$
(that is, we multiply all entries of this one row by $\lambda$, while leaving
all other entries of $A$ unchanged), then $\det A$ gets multiplied by
$\lambda$. \medskip

\textbf{(e)} If we add a row of $A$ to another row of $A$ (that is, we add
each entry of the former row to the corresponding entry of the latter), then
$\det A$ stays unchanged. \medskip

\textbf{(f)} Let $\lambda\in K$. If we add $\lambda$ times a row of $A$ to
another row of $A$ (that is, we add $\lambda$ times each entry of the former
row to the corresponding entry of the latter), then $\det A$ stays unchanged.
\medskip

\textbf{(g)} Let $B,C\in K^{n\times n}$ be two further $n\times n$-matrices.
Let $k\in\left[  n\right]  $. Assume that%
\[
\left(  \text{the }k\text{-th row of }C\right)  =\left(  \text{the }k\text{-th
row of }A\right)  +\left(  \text{the }k\text{-th row of }B\right)  ,
\]
whereas each $i\neq k$ satisfies%
\[
\left(  \text{the }i\text{-th row of }C\right)  =\left(  \text{the }i\text{-th
row of }A\right)  =\left(  \text{the }i\text{-th row of }B\right)  .
\]
Then,%
\[
\det C=\det A+\det B.
\]

\end{theorem}

\begin{example}
Let us see what Theorem \ref{thm.det.rowop} is saying in some particular cases
(specifically, for $3\times3$-matrices): \medskip

\textbf{(a)} One instance of Theorem \ref{thm.det.rowop} \textbf{(a)} is
\[
\det\left(
\begin{array}
[c]{ccc}%
a & b & c\\
a^{\prime\prime} & b^{\prime\prime} & c^{\prime\prime}\\
a^{\prime} & b^{\prime} & c^{\prime}%
\end{array}
\right)  =-\det\left(
\begin{array}
[c]{ccc}%
a & b & c\\
a^{\prime} & b^{\prime} & c^{\prime}\\
a^{\prime\prime} & b^{\prime\prime} & c^{\prime\prime}%
\end{array}
\right)  .
\]


\textbf{(b)} One instance of Theorem \ref{thm.det.rowop} \textbf{(b)} is
\[
\det\left(
\begin{array}
[c]{ccc}%
a & b & c\\
0 & 0 & 0\\
a^{\prime\prime} & b^{\prime\prime} & c^{\prime\prime}%
\end{array}
\right)  =0.
\]


\textbf{(c)} One instance of Theorem \ref{thm.det.rowop} \textbf{(c)} is
\[
\det\left(
\begin{array}
[c]{ccc}%
a & b & c\\
a^{\prime} & b^{\prime} & c^{\prime}\\
a & b & c
\end{array}
\right)  =0.
\]


\textbf{(d)} One instance of Theorem \ref{thm.det.rowop} \textbf{(d)} is
\[
\det\left(
\begin{array}
[c]{ccc}%
a & b & c\\
\lambda a^{\prime} & \lambda b^{\prime} & \lambda c^{\prime}\\
a^{\prime\prime} & b^{\prime\prime} & c^{\prime\prime}%
\end{array}
\right)  =\lambda\det\left(
\begin{array}
[c]{ccc}%
a & b & c\\
a^{\prime} & b^{\prime} & c^{\prime}\\
a^{\prime\prime} & b^{\prime\prime} & c^{\prime\prime}%
\end{array}
\right)  .
\]


\textbf{(e)} One instance of Theorem \ref{thm.det.rowop} \textbf{(e)} is
\[
\det\left(
\begin{array}
[c]{ccc}%
a & b & c\\
a^{\prime}+a^{\prime\prime} & b^{\prime}+b^{\prime\prime} & c^{\prime
}+c^{\prime\prime}\\
a^{\prime\prime} & b^{\prime\prime} & c^{\prime\prime}%
\end{array}
\right)  =\det\left(
\begin{array}
[c]{ccc}%
a & b & c\\
a^{\prime} & b^{\prime} & c^{\prime}\\
a^{\prime\prime} & b^{\prime\prime} & c^{\prime\prime}%
\end{array}
\right)  .
\]


\textbf{(f)} One instance of Theorem \ref{thm.det.rowop} \textbf{(f)} is
\[
\det\left(
\begin{array}
[c]{ccc}%
a & b & c\\
a^{\prime}+\lambda a^{\prime\prime} & b^{\prime}+\lambda b^{\prime\prime} &
c^{\prime}+\lambda c^{\prime\prime}\\
a^{\prime\prime} & b^{\prime\prime} & c^{\prime\prime}%
\end{array}
\right)  =\det\left(
\begin{array}
[c]{ccc}%
a & b & c\\
a^{\prime} & b^{\prime} & c^{\prime}\\
a^{\prime\prime} & b^{\prime\prime} & c^{\prime\prime}%
\end{array}
\right)  .
\]


\textbf{(g)} One instance of Theorem \ref{thm.det.rowop} \textbf{(g)} is%
\[
\det\underbrace{\left(
\begin{array}
[c]{ccc}%
a & b & c\\
d+d^{\prime} & e+e^{\prime} & f+f^{\prime}\\
g & h & i
\end{array}
\right)  }_{\text{this is }C}=\det\underbrace{\left(
\begin{array}
[c]{ccc}%
a & b & c\\
d & e & f\\
g & h & i
\end{array}
\right)  }_{\text{this is }A}+\det\underbrace{\left(
\begin{array}
[c]{ccc}%
a & b & c\\
d^{\prime} & e^{\prime} & f^{\prime}\\
g & h & i
\end{array}
\right)  }_{\text{this is }B}.
\]
(Specifically, this is the particular case of Theorem \ref{thm.det.rowop}
\textbf{(g)} for $n=3$ and $k=2$.)
\end{example}

Parts \textbf{(b)}, \textbf{(d)} and \textbf{(g)} of Theorem
\ref{thm.det.rowop} are commonly summarized under the mantle of
\textquotedblleft\emph{multilinearity of the determinant}\textquotedblright%
\ or \textquotedblleft\emph{linearity of the determinant in the }$k$\emph{-th
row}\textquotedblright. In fact, they say that (for any given $n\in\mathbb{N}$
and $k\in\left[  n\right]  $) if we hold all rows other than the $k$-th row of
an $n\times n$-matrix $A$ fixed, then $\det A$ depends $K$-linearly on the
$k$-th row of $A$.

\begin{proof}
[Proof of Theorem \ref{thm.det.rowop}.]\textbf{(a)} See \cite[Exercise 6.7
\textbf{(a)}]{detnotes}. This is also a particular case of \cite[Corollary
B.19]{Strick13}.

\textbf{(b)} See \cite[Exercise 6.7 \textbf{(c)}]{detnotes}. This is also
near-obvious from Definition \ref{def.det.det}.

\textbf{(c)} See \cite[Exercise 6.7 \textbf{(e)}]{detnotes} or \cite[\S 5.3.3,
property (iii)]{Laue-det} or \cite[2019-10-23 blackboard notes, Theorem
1.3.3]{19fla}.\footnote{\textit{Warning:} Several authors claim to give an
easy proof of part \textbf{(c)} using part \textbf{(a)}. This
\textquotedblleft proof\textquotedblright\ goes as follows: If $A$ has two
equal rows, then swapping these rows leaves $A$ unchanged, but (because of
Theorem \ref{thm.det.rowop}) flips the sign of $\det A$. Hence, in this case,
we have $\det A=-\det A$, so that $2\det A=0$ and therefore $\det A=0$, right?
Not so fast! In order to obtain $\det A=0$ from $2\det A=0$, we need the
element $2$ of $K$ to be invertible or at least be a non-zero-divisor (since
we have to divide by $2$). This is true when $K$ is one of the
\textquotedblleft high-school rings\textquotedblright\ $\mathbb{Z}$,
$\mathbb{Q}$, $\mathbb{R}$ and $\mathbb{C}$, but it is not true when $K$ is
the field $\mathbb{F}_{2}$ with $2$ elements (or, more generally, any field of
characteristic $2$). This slick argument can be salvaged, but in the form just
given it is incomplete.}

\textbf{(d)} See \cite[Exercise 6.7 \textbf{(g)}]{detnotes} or \cite[\S 5.3.3,
property (ii)]{Laue-det}. This is also a particular case of \cite[Corollary
B.19]{Strick13}.

\textbf{(f)} See \cite[Exercise 6.8 \textbf{(a)}]{detnotes}. This is also a
particular case of \cite[Corollary B.19]{Strick13}.

\textbf{(e)} This is the particular case of part \textbf{(f)} for $\lambda=1$.

\textbf{(g)} See \cite[Exercise 6.7 \textbf{(i)}]{detnotes} or \cite[\S 5.3.3,
property (i)]{Laue-det} or \cite[2019-10-30 blackboard notes, Theorem
1.2.3]{19fla}.
\end{proof}

\begin{theorem}
[Column operation properties]\label{thm.det.colop}Theorem \ref{thm.det.rowop}
also holds if we replace \textquotedblleft row\textquotedblright\ by
\textquotedblleft column\textquotedblright\ throughout it.
\end{theorem}

\begin{proof}
Theorem \ref{thm.det.transp} shows that the determinant of a matrix does not
change when we replace it by its transpose; however, the rows of this
transpose $A^{T}$ are the transposes of the columns of $A$. Thus, Theorem
\ref{thm.det.colop} follows by applying Theorem \ref{thm.det.rowop} to the
transposes of all the matrices involved. (See \cite[Exercises 6.7 and
6.8]{detnotes} for the details.)
\end{proof}

\begin{corollary}
\label{cor.det.sig-row-col}Let $n\in\mathbb{N}$. Let $A\in K^{n\times n}$ and
$\tau\in S_{n}$. Then,%
\begin{equation}
\det\left(  \left(  A_{\tau\left(  i\right)  ,j}\right)  _{1\leq i\leq
n,\ 1\leq j\leq n}\right)  =\left(  -1\right)  ^{\tau}\cdot\det A
\label{eq.cor.det.sig-row-col.row}%
\end{equation}
and%
\begin{equation}
\det\left(  \left(  A_{i,\tau\left(  j\right)  }\right)  _{1\leq i\leq
n,\ 1\leq j\leq n}\right)  =\left(  -1\right)  ^{\tau}\cdot\det A.
\label{eq.cor.det.sig-row-col.col}%
\end{equation}

\end{corollary}

In words: When we permute the rows or the columns of a matrix, its determinant
gets multiplied by the sign of the permutation.

\begin{proof}
[Proof of Corollary \ref{cor.det.sig-row-col}.]Let us first prove
(\ref{eq.cor.det.sig-row-col.col}).

The definition of $\det A$ yields%
\begin{align}
\det A  &  =\sum_{\sigma\in S_{n}}\left(  -1\right)  ^{\sigma}%
\underbrace{A_{1,\sigma\left(  1\right)  }A_{2,\sigma\left(  2\right)  }\cdots
A_{n,\sigma\left(  n\right)  }}_{=\prod_{i=1}^{n}A_{i,\sigma\left(  i\right)
}}\nonumber\\
&  =\sum_{\sigma\in S_{n}}\left(  -1\right)  ^{\sigma}\prod_{i=1}%
^{n}A_{i,\sigma\left(  i\right)  }. \label{pf.cor.det.sig-row-col.0}%
\end{align}


The definition of $\det\left(  \left(  A_{i,\tau\left(  j\right)  }\right)
_{1\leq i\leq n,\ 1\leq j\leq n}\right)  $ yields%
\begin{align}
\det\left(  \left(  A_{i,\tau\left(  j\right)  }\right)  _{1\leq i\leq
n,\ 1\leq j\leq n}\right)   &  =\sum_{\sigma\in S_{n}}\left(  -1\right)
^{\sigma}\underbrace{A_{1,\tau\left(  \sigma\left(  1\right)  \right)
}A_{2,\tau\left(  \sigma\left(  2\right)  \right)  }\cdots A_{n,\tau\left(
\sigma\left(  n\right)  \right)  }}_{=\prod_{i=1}^{n}A_{i,\tau\left(
\sigma\left(  i\right)  \right)  }}\nonumber\\
&  =\sum_{\sigma\in S_{n}}\left(  -1\right)  ^{\sigma}\prod_{i=1}^{n}%
A_{i,\tau\left(  \sigma\left(  i\right)  \right)  }.
\label{pf.cor.det.sig-row-col.1}%
\end{align}


However, $S_{n}$ is a group. Thus, the map%
\begin{align*}
S_{n}  &  \rightarrow S_{n},\\
\sigma &  \mapsto\tau^{-1}\sigma
\end{align*}
is a bijection. Hence, we can substitute $\tau^{-1}\sigma$ for $\sigma$ in the
sum $\sum_{\sigma\in S_{n}}\left(  -1\right)  ^{\sigma}\prod_{i=1}%
^{n}A_{i,\tau\left(  \sigma\left(  i\right)  \right)  }$. Thus, we obtain%
\begin{align*}
&  \sum_{\sigma\in S_{n}}\left(  -1\right)  ^{\sigma}\prod_{i=1}^{n}%
A_{i,\tau\left(  \sigma\left(  i\right)  \right)  }\\
&  =\sum_{\sigma\in S_{n}}\underbrace{\left(  -1\right)  ^{\tau^{-1}\sigma}%
}_{\substack{=\left(  -1\right)  ^{\tau^{-1}}\cdot\left(  -1\right)  ^{\sigma
}\\\text{(by Proposition \ref{prop.perm.sign.props} \textbf{(d)}%
,}\\\text{applied to }\tau^{-1}\text{ and }\sigma\text{ instead of }%
\sigma\text{ and }\tau\text{)}}}\prod_{i=1}^{n}\underbrace{A_{i,\tau\left(
\left(  \tau^{-1}\sigma\right)  \left(  i\right)  \right)  }}%
_{\substack{=A_{i,\sigma\left(  i\right)  }\\\text{(since }\tau\left(  \left(
\tau^{-1}\sigma\right)  \left(  i\right)  \right)  =\left(  \tau\tau
^{-1}\sigma\right)  \left(  i\right)  =\sigma\left(  i\right)
\\\text{(because }\tau\tau^{-1}\sigma=\sigma\text{))}}}\\
&  =\sum_{\sigma\in S_{n}}\left(  -1\right)  ^{\tau^{-1}}\cdot\left(
-1\right)  ^{\sigma}\prod_{i=1}^{n}A_{i,\sigma\left(  i\right)  }%
=\underbrace{\left(  -1\right)  ^{\tau^{-1}}}_{\substack{=\left(  -1\right)
^{\tau}\\\text{(by Proposition \ref{prop.perm.sign.props} \textbf{(f)}%
,}\\\text{applied to }\tau^{-1}\text{ instead of }\sigma\text{)}}%
}\cdot\underbrace{\sum_{\sigma\in S_{n}}\left(  -1\right)  ^{\sigma}%
\prod_{i=1}^{n}A_{i,\sigma\left(  i\right)  }}_{\substack{=\det A\\\text{(by
(\ref{pf.cor.det.sig-row-col.0}))}}}\\
&  =\left(  -1\right)  ^{\tau}\cdot\det A.
\end{align*}
In view of this, we can rewrite (\ref{pf.cor.det.sig-row-col.1}) as%
\[
\det\left(  \left(  A_{i,\tau\left(  j\right)  }\right)  _{1\leq i\leq
n,\ 1\leq j\leq n}\right)  =\left(  -1\right)  ^{\tau}\cdot\det A.
\]
This proves (\ref{eq.cor.det.sig-row-col.col}).

Now, we can easily derive (\ref{eq.cor.det.sig-row-col.row}) by using the
transpose. Indeed, applying (\ref{eq.cor.det.sig-row-col.col}) to $A^{T}$
instead of $A$, we obtain%
\begin{align}
\det\left(  \left(  \left(  A^{T}\right)  _{i,\tau\left(  j\right)  }\right)
_{1\leq i\leq n,\ 1\leq j\leq n}\right)   &  =\left(  -1\right)  ^{\tau}%
\cdot\underbrace{\det\left(  A^{T}\right)  }_{\substack{=\det A\\\text{(by
Theorem \ref{thm.det.transp})}}}\nonumber\\
&  =\left(  -1\right)  ^{\tau}\cdot\det A. \label{pf.cor.det.sig-row-col.5}%
\end{align}
However, each $\left(  i,j\right)  \in\left[  n\right]  ^{2}$ satisfies%
\[
\left(  A^{T}\right)  _{i,\tau\left(  j\right)  }=A_{\tau\left(  j\right)
,i}\ \ \ \ \ \ \ \ \ \ \left(  \text{by the definition of }A^{T}\right)  .
\]
Thus,%
\[
\left(  \left(  A^{T}\right)  _{i,\tau\left(  j\right)  }\right)  _{1\leq
i\leq n,\ 1\leq j\leq n}=\left(  A_{\tau\left(  j\right)  ,i}\right)  _{1\leq
i\leq n,\ 1\leq j\leq n}=\left(  \left(  A_{\tau\left(  i\right)  ,j}\right)
_{1\leq i\leq n,\ 1\leq j\leq n}\right)  ^{T}%
\]
(again by the definition of the transpose). Therefore,%
\begin{align*}
\det\left(  \left(  \left(  A^{T}\right)  _{i,\tau\left(  j\right)  }\right)
_{1\leq i\leq n,\ 1\leq j\leq n}\right)   &  =\det\left(  \left(  \left(
A_{\tau\left(  i\right)  ,j}\right)  _{1\leq i\leq n,\ 1\leq j\leq n}\right)
^{T}\right) \\
&  =\det\left(  \left(  A_{\tau\left(  i\right)  ,j}\right)  _{1\leq i\leq
n,\ 1\leq j\leq n}\right)
\end{align*}
(by Theorem \ref{thm.det.transp}, applied to $\left(  A_{\tau\left(  i\right)
,j}\right)  _{1\leq i\leq n,\ 1\leq j\leq n}$ instead of $A$). Hence,%
\[
\det\left(  \left(  A_{\tau\left(  i\right)  ,j}\right)  _{1\leq i\leq
n,\ 1\leq j\leq n}\right)  =\det\left(  \left(  \left(  A^{T}\right)
_{i,\tau\left(  j\right)  }\right)  _{1\leq i\leq n,\ 1\leq j\leq n}\right)
=\left(  -1\right)  ^{\tau}\cdot\det A
\]
(by (\ref{pf.cor.det.sig-row-col.5})). This proves
(\ref{eq.cor.det.sig-row-col.row}). Thus, the proof of Corollary
\ref{cor.det.sig-row-col} is complete.
\end{proof}

The following is probably the most remarkable property of determinants:

\begin{theorem}
[Multiplicativity of the determinant]\label{thm.det.detAB}Let $n\in\mathbb{N}%
$. Let $A,B\in K^{n\times n}$ be two $n\times n$-matrices. Then,%
\[
\det\left(  AB\right)  =\det A\cdot\det B.
\]

\end{theorem}

\begin{proof}
See \cite[Theorem B.17]{Strick13} or \cite[Theorem 6.23]{detnotes} or
\cite[\S 5]{Zeilbe} or \cite[Lemma 4.7.5 (1)]{Ford21} or \cite[Theorem
5.7]{Laue-det}. (Many of these proofs are at least partly combinatorial, but
the one in \cite[\S 5]{Zeilbe} is fully so, constructing quite explicitly a
sign-reversing involution.)
\end{proof}

As an application of these properties, let us reprove Proposition
\ref{prop.det.xiyj} and Proposition \ref{prop.det.xi+yj}:

\begin{proof}
[Second proof of Proposition \ref{prop.det.xiyj}.]Define two $n\times
n$-matrices%
\[
A:=\left(
\begin{array}
[c]{ccccc}%
x_{1} & 0 & 0 & \cdots & 0\\
x_{2} & 0 & 0 & \cdots & 0\\
x_{3} & 0 & 0 & \cdots & 0\\
\vdots & \vdots & \vdots & \ddots & \vdots\\
x_{n} & 0 & 0 & \cdots & 0
\end{array}
\right)  \ \ \ \ \ \ \ \ \ \ \text{and}\ \ \ \ \ \ \ \ \ \ B:=\left(
\begin{array}
[c]{ccccc}%
y_{1} & y_{2} & y_{3} & \cdots & y_{n}\\
0 & 0 & 0 & \cdots & 0\\
0 & 0 & 0 & \cdots & 0\\
\vdots & \vdots & \vdots & \ddots & \vdots\\
0 & 0 & 0 & \cdots & 0
\end{array}
\right)  .
\]
(Only the first column of $A$ and the first row of $B$ have any nonzero entries.)

Now,
\begin{align*}
\left(  x_{i}y_{j}\right)  _{1\leq i\leq n,\ 1\leq j\leq n}  &  =\left(
\begin{array}
[c]{cccc}%
x_{1}y_{1} & x_{1}y_{2} & \cdots & x_{1}y_{n}\\
x_{2}y_{1} & x_{2}y_{2} & \cdots & x_{2}y_{n}\\
\vdots & \vdots & \ddots & \vdots\\
x_{n}y_{1} & x_{n}y_{2} & \cdots & x_{n}y_{n}%
\end{array}
\right) \\
&  =\underbrace{\left(
\begin{array}
[c]{ccccc}%
x_{1} & 0 & 0 & \cdots & 0\\
x_{2} & 0 & 0 & \cdots & 0\\
x_{3} & 0 & 0 & \cdots & 0\\
\vdots & \vdots & \vdots & \ddots & \vdots\\
x_{n} & 0 & 0 & \cdots & 0
\end{array}
\right)  }_{=A}\underbrace{\left(
\begin{array}
[c]{ccccc}%
y_{1} & y_{2} & y_{3} & \cdots & y_{n}\\
0 & 0 & 0 & \cdots & 0\\
0 & 0 & 0 & \cdots & 0\\
\vdots & \vdots & \vdots & \ddots & \vdots\\
0 & 0 & 0 & \cdots & 0
\end{array}
\right)  }_{=B}=AB,
\end{align*}
so that%
\[
\det\left(  \left(  x_{i}y_{j}\right)  _{1\leq i\leq n,\ 1\leq j\leq
n}\right)  =\det\left(  AB\right)  =\det A\cdot\det B
\]
(by Theorem \ref{thm.det.detAB}). However, the matrix $A$ has a zero column
(since $n\geq2$), and thus satisfies $\det A=0$ (by Theorem
\ref{thm.det.colop} \textbf{(b)}\footnote{Of course, by \textquotedblleft
Theorem \ref{thm.det.colop} \textbf{(b)}\textquotedblright, we mean
\textquotedblleft the analogue of Theorem \ref{thm.det.rowop} \textbf{(b)} for
columns instead of rows\textquotedblright.}). Hence,%
\[
\det\left(  \left(  x_{i}y_{j}\right)  _{1\leq i\leq n,\ 1\leq j\leq
n}\right)  =\underbrace{\det A}_{=0}\cdot\,\det B=0.
\]
Thus, Proposition \ref{prop.det.xiyj} is proven again.
\end{proof}

\begin{proof}
[Second proof of Proposition \ref{prop.det.xi+yj}.]Define two $n\times
n$-matrices%
\[
A:=\left(
\begin{array}
[c]{ccccc}%
x_{1} & 1 & 0 & \cdots & 0\\
x_{2} & 1 & 0 & \cdots & 0\\
x_{3} & 1 & 0 & \cdots & 0\\
\vdots & \vdots & \vdots & \ddots & \vdots\\
x_{n} & 1 & 0 & \cdots & 0
\end{array}
\right)  \ \ \ \ \ \ \ \ \ \ \text{and}\ \ \ \ \ \ \ \ \ \ B:=\left(
\begin{array}
[c]{ccccc}%
1 & 1 & 1 & \cdots & 1\\
y_{1} & y_{2} & y_{3} & \cdots & y_{n}\\
0 & 0 & 0 & \cdots & 0\\
\vdots & \vdots & \vdots & \ddots & \vdots\\
0 & 0 & 0 & \cdots & 0
\end{array}
\right)  .
\]
(Only the first two columns of $A$ and the first two rows of $B$ have any
nonzero entries.)

Now,
\begin{align*}
\left(  x_{i}+y_{j}\right)  _{1\leq i\leq n,\ 1\leq j\leq n}  &  =\left(
\begin{array}
[c]{cccc}%
x_{1}+y_{1} & x_{1}+y_{2} & \cdots & x_{1}+y_{n}\\
x_{2}+y_{1} & x_{2}+y_{2} & \cdots & x_{2}+y_{n}\\
\vdots & \vdots & \ddots & \vdots\\
x_{n}+y_{1} & x_{n}+y_{2} & \cdots & x_{n}+y_{n}%
\end{array}
\right) \\
&  =\underbrace{\left(
\begin{array}
[c]{ccccc}%
x_{1} & 1 & 0 & \cdots & 0\\
x_{2} & 1 & 0 & \cdots & 0\\
x_{3} & 1 & 0 & \cdots & 0\\
\vdots & \vdots & \vdots & \ddots & \vdots\\
x_{n} & 1 & 0 & \cdots & 0
\end{array}
\right)  }_{=A}\underbrace{\left(
\begin{array}
[c]{ccccc}%
1 & 1 & 1 & \cdots & 1\\
y_{1} & y_{2} & y_{3} & \cdots & y_{n}\\
0 & 0 & 0 & \cdots & 0\\
\vdots & \vdots & \vdots & \ddots & \vdots\\
0 & 0 & 0 & \cdots & 0
\end{array}
\right)  }_{=B}=AB,
\end{align*}
so that%
\[
\det\left(  \left(  x_{i}+y_{j}\right)  _{1\leq i\leq n,\ 1\leq j\leq
n}\right)  =\det\left(  AB\right)  =\det A\cdot\det B
\]
(by Theorem \ref{thm.det.detAB}). However, the matrix $A$ has a zero column
(since $n\geq3$), and thus satisfies $\det A=0$ (by Theorem
\ref{thm.det.colop} \textbf{(b)}). Hence,%
\[
\det\left(  \left(  x_{i}+y_{j}\right)  _{1\leq i\leq n,\ 1\leq j\leq
n}\right)  =\underbrace{\det A}_{=0}\cdot\,\det B=0.
\]
Thus, Proposition \ref{prop.det.xi+yj} is proven again.
\end{proof}

The following fact follows equally easily from everything we know so far about determinants:

\begin{corollary}
\label{cor.det.scale-row-col}Let $n\in\mathbb{N}$. Let $A\in K^{n\times n}$
and $d_{1},d_{2},\ldots,d_{n}\in K$. Then,%
\begin{equation}
\det\left(  \left(  d_{i}A_{i,j}\right)  _{1\leq i\leq n,\ 1\leq j\leq
n}\right)  =d_{1}d_{2}\cdots d_{n}\cdot\det A
\label{eq.cor.det.scale-row-col.row}%
\end{equation}
and%
\begin{equation}
\det\left(  \left(  d_{j}A_{i,j}\right)  _{1\leq i\leq n,\ 1\leq j\leq
n}\right)  =d_{1}d_{2}\cdots d_{n}\cdot\det A.
\label{eq.cor.det.scale-row-col.col}%
\end{equation}

\end{corollary}

\begin{proof}
[First proof of Corollary \ref{cor.det.scale-row-col} (sketched).]The matrix
$\left(  d_{i}A_{i,j}\right)  _{1\leq i\leq n,\ 1\leq j\leq n}$ is obtained
from the matrix $A$ by multiplying the $1$-st row by $d_{1}$, multiplying the
$2$-nd row by $d_{2}$, multiplying the $3$-rd row by $d_{3}$, and so on.
Theorem \ref{thm.det.rowop} \textbf{(d)} (applied repeatedly -- once for each
row) shows that these multiplications result in the determinant of $A$ getting
multiplied by $d_{1},d_{2},\ldots,d_{n}$ (in succession). Hence,%
\[
\det\left(  \left(  d_{i}A_{i,j}\right)  _{1\leq i\leq n,\ 1\leq j\leq
n}\right)  =d_{1}d_{2}\cdots d_{n}\cdot\det A.
\]
Thus, (\ref{eq.cor.det.scale-row-col.row}) is proved. The proof of
(\ref{eq.cor.det.scale-row-col.col}) is analogous (using columns instead of
rows). Corollary \ref{cor.det.scale-row-col} is proven.
\end{proof}

\begin{proof}
[Second proof of Corollary \ref{cor.det.scale-row-col} (sketched).]Let $D$ be
the diagonal matrix
\[
\left(
\begin{array}
[c]{cccc}%
d_{1} & 0 & \cdots & 0\\
0 & d_{2} & \cdots & 0\\
\vdots & \vdots & \ddots & \vdots\\
0 & 0 & \cdots & d_{n}%
\end{array}
\right)  \in K^{n\times n}.
\]
Then, $D$ is upper-triangular; hence, Theorem \ref{thm.det.triang} shows that
its determinant is $\det D=d_{1}d_{2}\cdots d_{n}$.

However, it is easy to see that $\left(  d_{i}A_{i,j}\right)  _{1\leq i\leq
n,\ 1\leq j\leq n}=DA$. Hence,%
\begin{align*}
\det\left(  \left(  d_{i}A_{i,j}\right)  _{1\leq i\leq n,\ 1\leq j\leq
n}\right)   &  =\det\left(  DA\right)  =\det D\cdot\det
A\ \ \ \ \ \ \ \ \ \ \left(  \text{by Theorem \ref{thm.det.detAB}}\right) \\
&  =d_{1}d_{2}\cdots d_{n}\cdot\det A\ \ \ \ \ \ \ \ \ \ \left(  \text{since
}\det D=d_{1}d_{2}\cdots d_{n}\right)  .
\end{align*}
This proves (\ref{eq.cor.det.scale-row-col.row}). A similar argument using
$\left(  d_{j}A_{i,j}\right)  _{1\leq i\leq n,\ 1\leq j\leq n}=AD$ proves
(\ref{eq.cor.det.scale-row-col.col}). Thus, Corollary
\ref{cor.det.scale-row-col} is proven again.
\end{proof}

\begin{proof}
[Third proof of Corollary \ref{cor.det.scale-row-col}.]Let us now proceed
completely elementarily. The definition of a determinant yields%
\begin{equation}
\det A=\sum_{\sigma\in S_{n}}\left(  -1\right)  ^{\sigma}A_{1,\sigma\left(
1\right)  }A_{2,\sigma\left(  2\right)  }\cdots A_{n,\sigma\left(  n\right)  }
\label{pf.cor.det.scale-row-col.3rd.1}%
\end{equation}
and%
\begin{align*}
\det\left(  \left(  d_{i}A_{i,j}\right)  _{1\leq i\leq n,\ 1\leq j\leq
n}\right)   &  =\sum_{\sigma\in S_{n}}\left(  -1\right)  ^{\sigma
}\underbrace{\left(  d_{1}A_{1,\sigma\left(  1\right)  }\right)  \left(
d_{2}A_{2,\sigma\left(  2\right)  }\right)  \cdots\left(  d_{n}A_{n,\sigma
\left(  n\right)  }\right)  }_{=\left(  d_{1}d_{2}\cdots d_{n}\right)  \left(
A_{1,\sigma\left(  1\right)  }A_{2,\sigma\left(  2\right)  }\cdots
A_{n,\sigma\left(  n\right)  }\right)  }\\
&  =\sum_{\sigma\in S_{n}}\left(  -1\right)  ^{\sigma}\left(  d_{1}d_{2}\cdots
d_{n}\right)  \left(  A_{1,\sigma\left(  1\right)  }A_{2,\sigma\left(
2\right)  }\cdots A_{n,\sigma\left(  n\right)  }\right) \\
&  =d_{1}d_{2}\cdots d_{n}\cdot\underbrace{\sum_{\sigma\in S_{n}}\left(
-1\right)  ^{\sigma}A_{1,\sigma\left(  1\right)  }A_{2,\sigma\left(  2\right)
}\cdots A_{n,\sigma\left(  n\right)  }}_{\substack{=\det A\\\text{(by
(\ref{pf.cor.det.scale-row-col.3rd.1}))}}}\\
&  =d_{1}d_{2}\cdots d_{n}\cdot\det A
\end{align*}
and%
\begin{align*}
&  \det\left(  \left(  d_{j}A_{i,j}\right)  _{1\leq i\leq n,\ 1\leq j\leq
n}\right) \\
&  =\sum_{\sigma\in S_{n}}\left(  -1\right)  ^{\sigma}\underbrace{\left(
d_{\sigma\left(  1\right)  }A_{1,\sigma\left(  1\right)  }\right)  \left(
d_{\sigma\left(  2\right)  }A_{2,\sigma\left(  2\right)  }\right)
\cdots\left(  d_{\sigma\left(  n\right)  }A_{n,\sigma\left(  n\right)
}\right)  }_{=\left(  d_{\sigma\left(  1\right)  }d_{\sigma\left(  2\right)
}\cdots d_{\sigma\left(  n\right)  }\right)  \left(  A_{1,\sigma\left(
1\right)  }A_{2,\sigma\left(  2\right)  }\cdots A_{n,\sigma\left(  n\right)
}\right)  }\\
&  =\sum_{\sigma\in S_{n}}\left(  -1\right)  ^{\sigma}\underbrace{\left(
d_{\sigma\left(  1\right)  }d_{\sigma\left(  2\right)  }\cdots d_{\sigma
\left(  n\right)  }\right)  }_{\substack{=d_{1}d_{2}\cdots d_{n}\\\text{(since
}\sigma\text{ is a permutation of the set }\left[  n\right]  \text{)}}}\left(
A_{1,\sigma\left(  1\right)  }A_{2,\sigma\left(  2\right)  }\cdots
A_{n,\sigma\left(  n\right)  }\right) \\
&  =\sum_{\sigma\in S_{n}}\left(  -1\right)  ^{\sigma}\left(  d_{1}d_{2}\cdots
d_{n}\right)  \left(  A_{1,\sigma\left(  1\right)  }A_{2,\sigma\left(
2\right)  }\cdots A_{n,\sigma\left(  n\right)  }\right) \\
&  =d_{1}d_{2}\cdots d_{n}\cdot\det A\ \ \ \ \ \ \ \ \ \ \left(  \text{as we
have seen above}\right)  .
\end{align*}
Once again, this proves Corollary \ref{cor.det.scale-row-col}.
\end{proof}

